\chapter{Interfacce, classi astratte e generici in Java}
\label{chap:interfacce-generici-java}

Dopo l'ereditarietà e il polimorfismo, si approfondiscono i
comportamenti comuni agli oggetti, il controllo di accesso e le
interfacce. Si introducono inoltre le classi astratte e i tipi
generici (\cref{chap:realizzazione-classi-java,chap:ereditarieta-polimorfismo-java}).%
\footnote{Materiale di riferimento: S. Canazza,
\emph{Java: Interfacce},
\nolinkurl{2026_FI_15_canazza_Java_Interfacce.key.pdf}.}

\section{Obiettivi e principi della programmazione a oggetti}
\label{sec:fi-java-interfacce-progettazione}

La programmazione orientata agli oggetti organizza stato e
comportamento in componenti con responsabilità definite.
Incapsulamento, ereditarietà e polimorfismo, già introdotti nei
capitoli precedenti, contribuiscono a tre obiettivi:

\begin{table}[htbp]
\centering
\begin{tabular}{p{2.7cm}p{9.2cm}}
\toprule
Obiettivo & Significato \\
\midrule
Robustezza & Contenere le conseguenze di errori e condizioni inattese. \\
Adattabilità & Facilitare modifiche, estensioni e portabilità. \\
Riusabilità & Utilizzare componenti in programmi differenti senza duplicarne la logica. \\
\bottomrule
\end{tabular}
\caption{Obiettivi principali della progettazione orientata agli oggetti.}
\label{tab:fi-java-interfacce-obiettivi}
\end{table}

L'\emph{astrazione} individua gli aspetti rilevanti di un sistema;
l'\emph{information hiding} nasconde la rappresentazione interna
esponendo le operazioni necessarie; la \emph{modularità} separa il
software in unità funzionali. Questi principi favoriscono la
manutenzione, ma non eliminano i costi di gerarchie o dipendenze
progettate in modo inadeguato.

\section{La superclasse \texttt{Object}}
\label{sec:fi-java-interfacce-object}

Ogni classe Java deriva direttamente o indirettamente da
\texttt{Object} (\cref{def:fi-java-ereditarieta-object}).
Anche \texttt{println(Object)} può quindi ricevere un
riferimento a un oggetto e produrne una descrizione testuale
tramite \texttt{toString()}.

\subsection{Il metodo \texttt{toString()}}
\label{sec:fi-java-interfacce-tostring}

L'implementazione ereditata da \texttt{Object} restituisce
normalmente il nome della classe, il carattere \texttt{@} e un
codice hash in forma esadecimale. Tale codice non è
necessariamente l'indirizzo di memoria dell'oggetto. Una classe
può sovrascrivere il metodo per produrre una descrizione utile.

\begin{example}[Descrizione di \texttt{BankAccount}]
Si aggiunge il seguente metodo alla classe già utilizzata
nella \cref{sec:fi-java-ereditarieta-conti}:

\begin{unipdcode}[language=Java]
@Override
public String toString() {
    return getClass().getSimpleName()
        + "[balance=" + getBalance() + "]";
}
\end{unipdcode}


\begin{unipdcode}[language=Java]
BankAccount conto = new BankAccount(1000);
System.out.println(conto);
String descrizione = "Conto: " + conto;
\end{unipdcode}

\begin{unipdterminal}
BankAccount[balance=1000.0]
\end{unipdterminal}

La concatenazione antepone \texttt{"Conto: "} alla stessa
descrizione dell'oggetto.
\texttt{getClass()} restituisce un oggetto \texttt{Class}
relativo alla classe effettiva; \texttt{getSimpleName()} ne
restituisce il nome semplice, mentre \texttt{getName()} fornisce
il nome completo, comprensivo dell'eventuale package.
\end{example}

Una sottoclasse con stato aggiuntivo può sovrascrivere
\texttt{toString()}. Per riutilizzare la descrizione ereditata
può invocare \texttt{super.toString()}.
Le informazioni riservate non devono essere incluse
indiscriminatamente nella descrizione.

\subsection{Il metodo \texttt{equals()}}
\label{sec:fi-java-interfacce-equals}

L'implementazione di \texttt{Object.equals(Object)}
confronta l'identità dei riferimenti. Una sovrascrittura può
stabilire un'uguaglianza logica diversa da \texttt{==}
(\cref{def:fi-java-cicli-uguaglianza-oggetti}). Il parametro
deve avere tipo \texttt{Object}: dichiarare soltanto
\lstinline|equals(BankAccount other)| definirebbe un nuovo
sovraccarico e non sovrascriverebbe il metodo originale.

\begin{example}[Uguaglianza per classe e saldo]
I seguenti metodi, aggiunti a \texttt{BankAccount}, definiscono
l'uguaglianza quando classe effettiva e saldo coincidono:

\begin{unipdcode}[language=Java]
@Override
public boolean equals(Object other) {
    if (this == other) return true;
    if (other == null || getClass() != other.getClass())
        return false;
    BankAccount conto = (BankAccount) other;
    return Double.compare(getBalance(), conto.getBalance()) == 0;
}

@Override
public int hashCode() {
    return Double.hashCode(getBalance());
}
\end{unipdcode}


La sovrascrittura di \texttt{equals()} richiede coerenza con
\texttt{hashCode()}: due oggetti uguali devono avere lo stesso
codice hash. Due hash uguali non implicano invece l'uguaglianza
degli oggetti.
\end{example}

\section{Controllo di accesso e \texttt{final}}
\label{sec:fi-java-interfacce-accesso}

Ai modificatori \texttt{public} e \texttt{private} già trattati
nella \cref{sec:fi-java-classi-incapsulamento} si aggiungono
\texttt{protected} e l'accesso di package:

\clearpage

\begin{table}[htbp]
\centering
\begin{tabular}{p{3.0cm}p{9.0cm}}
\toprule
Modificatore & Visibilità dei membri \\
\midrule
\texttt{public} & Da ogni classe che possa accedere alla classe dichiarante. \\
\texttt{protected} & Dalle classi del medesimo package e dalle sottoclassi, secondo le regole di Java. \\
Nessuno (package) & Dalle sole classi dello stesso package. \\
\texttt{private} & All'interno della classe dichiarante. \\
\bottomrule
\end{tabular}
\caption{Controllo di accesso ai membri di una classe.}
\label{tab:fi-java-interfacce-accesso}
\end{table}

Una sottoclasse posta in un altro package non può usare un
membro \texttt{protected} tramite un riferimento arbitrario alla
superclasse. Per una classe di primo livello sono consentiti solo
\texttt{public} e l'accesso di package. Esporre direttamente
campi \texttt{protected} può creare dipendenze dalla rappresentazione
interna e rendere più difficili le modifiche successive.

\texttt{final} produce effetti diversi secondo la dichiarazione:

\begin{table}[htbp]
\centering
\begin{tabular}{p{2.6cm}p{6.5cm}p{3.0cm}}
\toprule
Dichiarazione & Effetto & Esempio \\
\midrule
Variabile & Impedisce successive assegnazioni. & \texttt{final int N = 4;} \\
Metodo & Impedisce la sovrascrittura. & \texttt{final void run()} \\
Classe & Impedisce la creazione di sottoclassi. & \texttt{final class C} \\
\bottomrule
\end{tabular}
\caption{Utilizzi principali del modificatore \texttt{final}.}
\label{tab:fi-java-interfacce-final}
\end{table}

Un riferimento \texttt{final} non può essere riassegnato, ma
l'oggetto referenziato può rimanere modificabile. La classe
\texttt{String} è \texttt{final} e i suoi oggetti sono immutabili:
si tratta di proprietà distinte.

\section{Interfacce e astrazione}
\label{sec:fi-java-interfacce-astrazione}

\begin{definition}[Interfaccia Java]
\label{def:fi-java-interfaccia}
Un'\emph{interfaccia} definisce un contratto di comportamento
comune a più classi. Si dichiara con \texttt{interface}; una
classe indica i contratti realizzati mediante \texttt{implements}.
\end{definition}

Un'interfaccia non ha costruttori o campi di istanza. I suoi
metodi astratti sono implicitamente \texttt{public}, mentre le
costanti sono \texttt{public static final}. Può contenere
metodi \texttt{default} e \texttt{static} con implementazione.

\subsection{Dichiarazione, implementazione e metodi \texttt{default}}
\label{sec:fi-java-interfacce-implementazione}

\begin{example}[Interfaccia \texttt{Measurable}]
L'interfaccia specifica una misura e fornisce un confronto
predefinito e una funzione statica di aggregazione:

\noindent
\begin{minipage}{\linewidth}
\begin{unipdcode}[language=Java]
public interface Measurable {
    double getMeasure();

    default boolean smallerThan(Measurable other) {
        return getMeasure() < other.getMeasure();
    }

    static double average(Measurable[] values) {
        if (values.length == 0)
            throw new IllegalArgumentException("Array vuoto");
        double sum = 0;
        for (Measurable item : values)
            sum += item.getMeasure();
        return sum / values.length;
    }
}
\end{unipdcode}

\end{minipage}

Una classe che implementa \texttt{Measurable} deve fornire
\texttt{getMeasure()}, mentre può ereditare \texttt{smallerThan()}:

\begin{unipdcode}[language=Java]
public class MeasuredItem implements Measurable {
    private final double measure;

    public MeasuredItem(double measure) {
        this.measure = measure;
    }

    @Override
    public double getMeasure() {
        return measure;
    }
}
\end{unipdcode}


\begin{unipdcode}[language=Java]
Measurable[] items = {new MeasuredItem(2), new MeasuredItem(4)};
double media = Measurable.average(items); // 3.0
boolean minore = items[0].smallerThan(items[1]); // true
\end{unipdcode}

\texttt{Measurable.average(array)} si invoca attraverso il
nome dell'interfaccia; richiede un array non vuoto e non nullo,
con elementi non nulli.
\end{example}

Una classe può implementare più interfacce, pur estendendo una
sola classe. Un'interfaccia può estendere più interfacce.
L'implementazione di un contratto è distinta dall'ereditarietà
della rappresentazione di una classe.

\subsection{Riferimenti, conversioni e polimorfismo}
\label{sec:fi-java-interfacce-riferimenti}

Un riferimento può avere come tipo dichiarato un'interfaccia;
non è però possibile costruire un oggetto mediante
\lstinline|new Measurable()|. L'assegnazione dal tipo di una
classe che la implementa al tipo dell'interfaccia è implicita:

\begin{unipdcode}[language=Java]
Measurable misura = new MeasuredItem(12);
System.out.println(misura.getMeasure());
\end{unipdcode}

Tramite \texttt{misura} si possono invocare i metodi disponibili
dal contratto, non i metodi specifici di \texttt{MeasuredItem}.
La conversione inversa richiede un cast valido per l'oggetto
effettivo; un controllo con \texttt{instanceof} può prevenirne
l'errore (\cref{sec:fi-java-ereditarieta-instanceof}).

Classi differenti possono fornire implementazioni diverse di
\texttt{getMeasure()}. La chiamata attraverso un riferimento
\texttt{Measurable} seleziona durante l'esecuzione il metodo
della classe effettiva, secondo il polimorfismo già studiato
nella \cref{sec:fi-java-ereditarieta-polimorfismo}.

\section{L'interfaccia \texttt{Comparable}}
\label{sec:fi-java-interfacce-comparable}

\begin{definition}[Ordinamento naturale]
\label{def:fi-java-interfacce-comparable}
\texttt{Comparable<T>} è l'interfaccia di \texttt{java.lang}
che dichiara \texttt{int compareTo(T other)} per stabilire
l'ordinamento naturale fra oggetti confrontabili. Il valore
restituito è negativo, nullo o positivo se l'oggetto corrente
precede, equivale o segue l'argomento nell'ordine.
\end{definition}

Non è corretto verificare \lstinline|compareTo(...) == -1|:
il risultato negativo non è necessariamente \(-1\).
Il criterio deve rispettare le proprietà di un ordinamento
coerente, tra cui riflessività e transitività della relazione
non stretta e compatibilità dei segni dei confronti inversi.
È raccomandabile, ma non obbligatorio, che
\texttt{compareTo(other) == 0} coincida con \texttt{equals(other)}.

\begin{example}[Confronto dei conti per saldo]
Alla dichiarazione di \texttt{BankAccount} si aggiunge
\lstinline|implements Comparable<BankAccount>| e si introduce:

\begin{unipdcode}[language=Java]
@Override
public int compareTo(BankAccount other) {
    return Double.compare(getBalance(), other.getBalance());
}
\end{unipdcode}


\texttt{Double.compare()} confronta anche i valori particolari
di \texttt{double}. Questo criterio usa esclusivamente il saldo:
può quindi considerare equivalenti nell'ordine oggetti che
\texttt{equals()} distingue in base alla classe effettiva.
\end{example}

\texttt{String} implementa \texttt{Comparable<String>} e
utilizza l'ordinamento lessicografico introdotto nella
\cref{sec:fi-java-ereditarieta-ordinamento-oggetti}.

\section{La classe \texttt{ArrayAlgs} con \texttt{Comparable}}
\label{sec:fi-java-interfacce-arrayalgs}

Gli algoritmi della \cref{chap:ricerca-ordinamento-java} restano
validi per oggetti confrontabili. Cambiano il tipo dell'array
e il confronto, non la logica dell'algoritmo:

\begin{table}[htbp]
\centering
\begin{tabular}{p{4.0cm}p{4.4cm}p{3.8cm}}
\toprule
Relazione & Array di \texttt{int} & Array di \texttt{Comparable} \\
\midrule
Precedenza & \lstinline|a < b| & \lstinline|a.compareTo(b) < 0| \\
Posteriorità & \lstinline|a > b| & \lstinline|a.compareTo(b) > 0| \\
Equivalenza nell'ordine & \lstinline|a == b| & \lstinline|a.compareTo(b) == 0| \\
\bottomrule
\end{tabular}
\caption{Adattamento dei confronti numerici agli oggetti.}
\label{tab:fi-java-interfacce-confronti}
\end{table}

\subsection{Ordinamento per selezione}
\label{sec:fi-java-interfacce-selection}

La versione seguente mostra solo il metodo e i suoi ausiliari:

\noindent
\begin{minipage}{\linewidth}
\begin{unipdcode}[language=Java]
public static void selectionSort(Comparable[] v) {
    for (int i = 0; i < v.length - 1; i++) {
        int min = findMinPos(v, i);
        swap(v, i, min);
    }
}

private static int findMinPos(Comparable[] v, int from) {
    int min = from;
    for (int i = from + 1; i < v.length; i++)
        if (v[i].compareTo(v[min]) < 0) min = i;
    return min;
}

private static void swap(Comparable[] v, int i, int j) {
    Comparable temp = v[i];
    v[i] = v[j];
    v[j] = temp;
}
\end{unipdcode}

\end{minipage}

\begin{example}[Riutilizzo del metodo su tipi diversi]
La stessa implementazione ordina un array di stringhe o di conti
che implementino \texttt{Comparable}:

\begin{unipdcode}[language=Java]
String[] parole = {"Zoe", "Anna", "Luca"};
ArrayAlgs.selectionSort(parole);

BankAccount[] conti = {
    new BankAccount(200), new BankAccount(50)
};
ArrayAlgs.selectionSort(conti);
\end{unipdcode}

Le stringhe risultano in ordine lessicografico e i conti in
ordine crescente di saldo.
\end{example}

\texttt{findMinPos()} cerca il minimo con \texttt{compareTo()},
mentre \texttt{swap()} scambia i riferimenti. Assumendo
confronti a costo costante, il tempo resta \(\Theta(n^2)\)
e lo spazio ausiliario \(\Theta(1)\). Se confrontare gli
elementi è costoso, il tempo totale deve includere tale costo.

\subsection{Adattamento degli altri algoritmi}
\label{sec:fi-java-interfacce-altri-algoritmi}

In \texttt{insertionSort()}, la condizione di spostamento diventa
\lstinline|temp.compareTo(v[j-1]) < 0|. In \texttt{mergeSort()},
la fusione può confrontare
\lstinline|left[i].compareTo(right[j]) <= 0|, conservando la
stabilità a parità di confronto. In \texttt{binarySearch()},
si esamina il segno di \lstinline|v[mid].compareTo(value)|,
conservando l'ordinamento iniziale come precondizione.

L'uso non parametrizzato di \texttt{Comparable[]} presuppone
che gli elementi siano non nulli e reciprocamente confrontabili:
un confronto incompatibile può generare
\texttt{ClassCastException}. I tipi generici possono rendere
questo vincolo più esplicito al compilatore.

\section{Classi e metodi astratti}
\label{sec:fi-java-interfacce-astratte}

\begin{definition}[Classe e metodo astratti]
\label{def:fi-java-interfacce-astratte}
Una \emph{classe astratta}, dichiarata con \texttt{abstract},
non è direttamente istanziabile e può contenere metodi
implementati e metodi astratti. Un \emph{metodo astratto}
dichiara un'operazione senza fornirne il corpo.
\end{definition}

\begin{example}[Calcolo dell'area]
\texttt{Poligono} specifica il calcolo dell'area e realizza il
confronto in base a tale misura:

\noindent
\begin{minipage}{\linewidth}
\begin{unipdcode}[language=Java]
abstract class Poligono implements Comparable<Poligono> {
    public abstract double area();

    @Override
    public int compareTo(Poligono other) {
        return Double.compare(area(), other.area());
    }
}

class Rettangolo extends Poligono {
    private final double base, altezza;

    Rettangolo(double base, double altezza) {
        this.base = base;
        this.altezza = altezza;
    }

    @Override
    public double area() {
        return base * altezza;
    }
}
\end{unipdcode}

\end{minipage}

\texttt{Poligono} non può essere istanziata; la sottoclasse
concreta \texttt{Rettangolo} implementa \texttt{area()}.
Si assumono lati non negativi.
\end{example}

Una classe astratta può estendere un'altra classe astratta senza
realizzare subito tutti i metodi astratti; una sottoclasse
concreta deve invece implementarli. A differenza di
un'interfaccia, una classe astratta può avere campi di
istanza e costruttori.

\section{Introduzione alle classi generiche}
\label{sec:fi-java-interfacce-generici}

\begin{definition}[Parametrizzazione dei tipi]
\label{def:fi-java-interfacce-generici}
Una \emph{classe generica} dichiara uno o più parametri di tipo
tra parentesi angolari. Un utilizzo della classe specifica un
tipo concreto, controllato durante la compilazione.
\end{definition}

\lstinline|class Box<T>| dichiara un parametro \texttt{T};
\lstinline|Box<String>| e \lstinline|Box<Integer>| sono due
tipi parametrizzati. Le lettere \texttt{T}, \texttt{E},
\texttt{K}, \texttt{V}, \texttt{N} sono convenzioni rispettivamente
per tipo, elemento, chiave, valore e tipo numerico. I parametri
di tipo richiedono tipi riferimento, quindi si usa
\texttt{Integer} al posto di \texttt{int}.

\begin{example}[Contenitore \texttt{Box<T>}]
La classe conserva un oggetto del tipo indicato dal parametro:

\noindent
\begin{minipage}{\linewidth}
\begin{unipdcode}[language=Java]
public class Box<T> {
    private T content;

    public void set(T value) {
        content = value;
    }

    public T get() {
        return content;
    }
}
\end{unipdcode}

\end{minipage}

\begin{unipdcode}[language=Java]
Box<String> parole = new Box<>();
parole.set("ciao");
String s = parole.get();

Box<Integer> numeri = new Box<>();
numeri.set(42);
Integer n = numeri.get();
\end{unipdcode}

La chiamata \texttt{parole.set(42)} non è valida: un
\texttt{Integer} non soddisfa il parametro \texttt{String}.
Il recupero tramite \texttt{get()} non richiede cast.
\end{example}

\subsection{Metodi generici}
\label{sec:fi-java-interfacce-metodi-generici}

Un metodo generico dichiara i parametri di tipo prima del tipo
di ritorno:

\begin{unipdcode}[language=Java]
public static <T> void printArray(T[] array) {
    for (T element : array)
        System.out.print(element + " ");
    System.out.println();
}
\end{unipdcode}


\begin{unipdcode}[language=Java]
printArray(new String[] {"A", "B"});
printArray(new Integer[] {1, 2, 3});
\end{unipdcode}

Il compilatore deduce \texttt{T} dall'argomento. La stessa
implementazione può quindi operare su array di tipi riferimento
diversi, conservando i controlli statici sui tipi.
